\begin{lemma}
The exact final colored graph \(G(\omega)\) defined in
\(\noderef{TwoBiteFinalGraph}\) contains no red triangle, no blue triangle, and
no mixed triangle with two edges of one color and one edge of the other.  Hence
the simple graph obtained by forgetting colors and merging double edges is
triangle-free.
\end{lemma}
\begin{proof}
Write \(G(\omega)=(G_R,G_B,G_0)\), where \(G_R\) is the final red graph,
\(G_B\) is the final blue graph, and \(G_0\) is the simple shadow graph.  The
first assertion, that \(G_R\) has no red triangle, is exactly
\(\noderef{TwoBiteFinalGraphNoRedTriangle}\).  The second assertion, that
\(G_B\) has no blue triangle, is exactly
\(\noderef{TwoBiteFinalGraphNoBlueTriangle}\).

The third assertion is exactly
\(\noderef{TwoBiteFinalGraphNoRedRedBlueTriangle}\): no final triangle has two
red edges incident to a common center and the opposite edge blue.  The fourth
assertion is exactly \(\noderef{TwoBiteFinalGraphNoBlueBlueRedTriangle}\): no
final triangle has two blue edges incident to a common center and the opposite
edge red.

Finally, \(\noderef{TwoBiteFinalGraphShadowTriangleFree}\) applies these four
colored exclusions to the shadow graph.  Every shadow edge is witnessed by a
surviving colored edge, so a shadow triangle would choose three surviving
colored edges on the same three vertices.  The chosen colors are either all
red, all blue, two red and one blue, or two blue and one red, and those four
possibilities are precisely the four exclusions above.  Thus \(G_0\) has no
triangle on three distinct vertices.  Combining the five conclusions gives the
claimed conjunction.
\end{proof}
